\documentclass[11pt]{article}
\usepackage{fltcourse}

\title{Level 5: Hardly ramified representations}
\author{Kevin Buzzard}
\date{}

\begin{document}
\maketitle

\noindent
The object of this game is to prove Fermat's Last Theorem (FLT). Each level's
boss is a theorem of the form ``Statement $B_n$ implies FLT''; beating the
boss means proving that $B_n$ implies the boss statement $B_{n-1}$ of the
previous level, so that the implications chain down to $B_1={}$FLT. We are
allowed to assume, without proof, any theorem published in the 1980s or
earlier, provided we state it cleanly. We are also free to use, without
comment, any construction which is a standard part of the mathematical
literature (and, in this project, is formalized in Lean or in its mathematics
library \emph{mathlib}). What we may \emph{not} assume are the modularity and
automorphy results of the 1990s and later; those we must prove.

\medskip\noindent
By the end of the previous four levels, FLT has been reduced to Statement
$B_4$. Recall the ingredients.
A \emph{Frey package} is a prime $p\geq 5$ together with three nonzero,
pairwise coprime integers $a,b,c$ with $a\equiv 3\pmod 4$, $b$ even, and
$a^p+b^p=c^p$. Its \emph{Frey curve} is the elliptic curve $E/\Q$ with
Weierstrass equation
\[
  Y^2 = X\,(X-a^p)(X+b^p).
\]
Fixing a separable closure $\Qbar$ of $\Q$, the $p$-torsion $E(\Qbar)[p]$ is a
two-dimensional $\F_p$-vector space carrying a continuous action of
$\GQ$, the \emph{mod $p$ representation}
$\rhobar_{E,p}\colon\GQ\to\Aut_{\F_p}\!\bigl(E(\Qbar)[p]\bigr)$.

\begin{Bn}{B_4}
For every Frey package $(a,b,c,p)$, every separable closure $\Qbar/\Q$, and the
associated Frey curve $E$, the representation $\rhobar_{E,p}$ is reducible.
\end{Bn}

\noindent
(Recall that a representation of a group on a vector space is \emph{reducible}
if it has a nonzero proper stable subspace, and \emph{irreducible} otherwise.)
The boss of the previous level was the theorem that $B_4$ implies FLT.

\medskip\noindent
In this level we do not prove anything hard. Instead we isolate exactly which
features of $\rhobar_{E,p}$ are responsible for its reducibility, package them
into a self-contained condition on a Galois representation --- being
\emph{hardly ramified} --- and reduce $B_4$ to a statement, $B_5$, in which no
elliptic curve appears. The whole of the rest of the course is then devoted to
proving $B_5$.

\section{Recollections}

We recall, briefly, the standard notions on which the definition rests. All of
them are part of the standard literature.

A \emph{coefficient ring} is a commutative local ring $R$ equipped with a
topology making it a compact Hausdorff topological ring. Its maximal ideal
$\m$ is then open, and the residue field $R/\m$ is finite; the characteristic
of $R/\m$ is a prime, the \emph{residue characteristic} of $R$. Examples: any
finite local ring, such as $\Z/\ell^n\Z$; the $\ell$-adic integers $\Zl$; and
the ring of integers $\calO_K$ of a finite extension $K/\Ql$, with its
$\m$-adic topology. Every coefficient ring is the projective limit of its
finite discrete quotients $R/J$ ($J$ ranging over the open ideals). We shall
use one standard fact: \emph{a coefficient ring of residue characteristic
$\ell$ carries a unique continuous ring homomorphism $\Zl\to R$}, so is
canonically a $\Zl$-algebra.

Let $R$ be a coefficient ring and $V$ a free $R$-module of finite rank. Then
$\Aut_R(V)$ is a topological group. By an \emph{$R$-representation} of $\GQ$ we
mean a continuous homomorphism $\rho\colon\GQ\to\Aut_R(V)$, where $\GQ$ carries
its profinite topology. Choosing an $R$-basis of $V$ identifies $\Aut_R(V)$
with $\GL_n(R)$, but the topology --- and hence the notion of continuity ---
does not depend on the choice.

For each prime $p$, an embedding $\Qbar\hookrightarrow\Qpbar$ realises the local
group $\GQp=\Gal(\Qpbar/\Qp)$ as a \emph{decomposition subgroup} $D_p\subseteq
\GQ$, with \emph{inertia subgroup} $I_p\subseteq D_p$. A different embedding
conjugates $D_p$ and $I_p$ by an element of $\GQ$. A \emph{Frobenius element}
$\Frob_p\in D_p$ is any element mapping to the arithmetic Frobenius
$x\mapsto x^p$ of the residue field; it is well defined modulo $I_p$. An
$R$-representation $\rho$ is \emph{unramified at $p$} if $I_p\subseteq\ker\rho$
for one (equivalently every) choice of $I_p$; then $\rho(\Frob_p)$ is
well-defined up to conjugacy, so its trace and determinant are honest
invariants.

The \emph{$\ell$-adic cyclotomic character}
$\chi=\chi_\ell\colon\GQ\to\Zl^\times$ is characterised by
$\sigma(\zeta)=\zeta^{\chi(\sigma)}$ for every $\ell$-power root of unity
$\zeta$; its reduction modulo $\ell$ is $\chibar\colon\GQ\to\Fl^\times$. It is
unramified outside $\ell$, with $\chi(\Frob_p)=p$ for $p\neq\ell$. If $R$ is a
coefficient ring of residue characteristic $\ell$, we regard $\chi$ as a
character $\GQ\to R^\times$ via the canonical map $\Zl\to R$.

Finally we recall \emph{flatness}, the one notion which is genuinely subtle,
because it is the arithmetic-geometric input that makes everything work. Let
$K/\Ql$ be a finite extension. A finite abelian group $M$ with a continuous
action of $\GK=\Gal(\Kbar/K)$ is \emph{flat} if it arises, $\GK$-equivariantly,
as the group $A(\Kbar)$ of $\Kbar$-points of a finite flat commutative group
scheme $A$ over $\calO_K$ --- equivalently, if there is a finite free
$\calO_K$-Hopf algebra whose $\Kbar$-points, with their Galois action, recover
$M$. A continuous $R$-representation $\rho\colon\GK\to\Aut_R(V)$ (with $R$ a
coefficient ring of residue characteristic $\ell$, $V$ free) is \emph{flat} if
$V/JV$ is flat in the previous sense for every open ideal $J\subseteq R$. We
shall use the following stability properties, all of which were known in the
1980s and go back to Raynaud's 1974 paper \emph{Sch\'emas en groupes de type
$(p,\dots,p)$}: a subrepresentation or quotient of a flat representation is
flat, and a finite direct sum of flat representations is flat.

For a \emph{global} representation $\rho\colon\GQ\to\Aut_R(V)$ with $R$ of
residue characteristic $\ell$, we say $\rho$ is \emph{flat at $\ell$} if its
restriction to a decomposition group $D_\ell\cong\GQl$ (via any choice of
$\Qbar\hookrightarrow\Qlbar$) is flat. Different choices give conjugate
representations, so the condition is independent of the choice.

\begin{remark}
Flatness, not being unramified, is the correct ``good behaviour'' condition at
the residue characteristic. The mod $\ell$ cyclotomic character
$\chibar\colon\GQl\to\Fl^\times$ is ramified (tamely so, factoring through the
totally ramified extension $\Ql(\zeta_\ell)/\Ql$), yet it is flat: it is cut
out by the group scheme $\mu_\ell$ of $\ell$-th roots of unity. Its $\ell$-adic
lift $\chi\colon\GQl\to\Zl^\times$ is even wildly ramified, and still flat. So
even the simplest interesting characters are ramified but flat.
\end{remark}

\section{Hardly ramified representations}

We can now give the central definition of the level. The adjective ``hardly
ramified'' is bespoke to this course: it names precisely the collection of
local conditions satisfied by the mod $p$ representation of a Frey curve.

\begin{definition}\label{def:hr}
Let $R$ be a coefficient ring of \emph{odd} residue characteristic $\ell$, let
$V$ be a free $R$-module of rank $2$, and let $\rho\colon\GQ\to\Aut_R(V)$ be a
continuous $R$-representation. We say $\rho$ is \emph{hardly ramified} if it
satisfies the following four conditions.
\begin{enumerate}
\item[\textup{(HR1)}] \emph{(cyclotomic determinant)} $\det\rho=\chi$, the
$\ell$-adic cyclotomic character, viewed in $R^\times$ via $\Zl\to R$.
\item[\textup{(HR2)}] \emph{(ramification)} $\rho$ is unramified outside $2$
and $\ell$; that is, $\rho$ is unramified at every prime $p\neq 2,\ell$.
\item[\textup{(HR3)}] \emph{(flat at $\ell$)} $\rho$ is flat at $\ell$.
\item[\textup{(HR4)}] \emph{(hardly ramified at $2$)} There is a
$\Gal(\Qbar_2/\Q_2)$-stable $R$-line $W\subseteq V$ with free rank one quotient
$V/W$, on which $\Gal(\Qbar_2/\Q_2)$ acts through an \emph{unramified}
character $\psi$ with $\psi^2=1$.
\end{enumerate}
\end{definition}

\noindent
Condition (HR4) can be phrased matricially: for a suitable basis,
\[
  \rho|_{\Gal(\Qbar_2/\Q_2)}\;\cong\;
  \begin{pmatrix}\psi\chi & * \\ 0 & \psi\end{pmatrix},
\]
with $\psi$ unramified, $\psi^2=1$, and $*$ unconstrained. (The upper-left
entry is forced to be $\psi\chi$ by \textup{(HR1)}, since $\det=\chi$ and $\chi$
is unramified at $2$.) It is the one place where we permit a controlled amount
of ramification --- hence the name.

\begin{example}\label{ex:1pluschi}
The direct sum $\mathbf 1\oplus\chi$ of the trivial character and the
cyclotomic character is hardly ramified over any coefficient ring of odd
residue characteristic. Indeed $\det(\mathbf 1\oplus\chi)=\chi$; it is
unramified away from $\ell$; both $\mathbf 1$ and $\chi$ are flat at $\ell$ (by
$\Z/\ell$ and $\mu_\ell$), so their sum is; and at $2$ one takes $W$ the
$\chi$-line and $\psi=\mathbf 1$. This is the representation the whole game is
trying to trap.
\end{example}

The boss statement of this level asserts that, for coefficients $\F_\ell$ with
$\ell\geq5$, Example~\ref{ex:1pluschi} is essentially the only possibility: no
hardly ramified representation can be irreducible.

\begin{Bn}{B_5}
Let $\ell\geq5$ be prime, let $V$ be a two-dimensional $\F_\ell$-vector space,
and let $\rho\colon\GQ\to\Aut_{\F_\ell}(V)$ be a continuous hardly ramified
representation. Then $\rho$ is reducible.
\end{Bn}

\section{The boss: $B_5$ implies FLT}

The link between $B_5$ and the Frey curve is a theorem of Serre from the
1980s, which we assume.

\begin{theorem}[Serre, 1987]\label{thm:frey-hr}
Let $(a,b,c,p)$ be a Frey package with Frey curve $E$, and let $\Qbar$ be any
separable closure of $\Q$. Then the mod $p$ representation
$\rhobar_{E,p}\colon\GQ\to\Aut_{\F_p}(E(\Qbar)[p])$ is hardly ramified.
\end{theorem}

This is proved in Serre's paper \emph{Sur les repr\'esentations modulaires de
degr\'e $2$ de $\Gal(\overline{\Q}/\Q)$} (Duke Math.\ J.\ \textbf{54}, 1987),
whose results were all available in the 1980s. We record the shape of the
argument, but as it is pre-1990 we give no proof.

\begin{itemize}
\item \emph{(HR1)}: the Weil pairing on $E[p]$ is an alternating perfect
$\GQ$-equivariant pairing $E[p]\times E[p]\to\mu_p$, and its existence forces
$\det\rhobar_{E,p}=\chi$.
\item \emph{(HR2), (HR4)}: pairwise coprimality of $a,b,c$ makes $E$ semistable
at every odd prime dividing $abc$, and the theory of the Tate curve shows the
representation is unramified there. At $2$ the same theory, together with the
normalisations $a\equiv3\pmod4$ and $b$ even (which put $E$ into a minimal
model with multiplicative reduction at $2$), yields precisely the shape of
\textup{(HR4)}.
\item \emph{(HR3)}: if $p\nmid abc$ then $E$ has good reduction at $p$ and
flatness is standard; if $p\mid abc$ the Tate-curve description again applies,
and one checks flatness using that $v_p(j_E)$ is divisible by~$p$.
\end{itemize}

\begin{theorem}[Boss of Level 5]\label{thm:boss5}
Statement $B_5$ implies Fermat's Last Theorem, modulo results known in the
1980s.
\end{theorem}

\begin{proof}
By the boss theorem of the previous level, it suffices to prove that $B_5$
implies Statement $B_4$. Let $(a,b,c,p)$ be a Frey package, $\Qbar$ a separable
closure of $\Q$, $E$ the Frey curve, and
$\rhobar_{E,p}\colon\GQ\to\Aut_{\F_p}(E(\Qbar)[p])$ the mod $p$ representation.
Since a Frey package has $p\geq5$, in particular $p$ is odd, so the
coefficient ring $\F_p$ has odd residue characteristic and
Definition~\ref{def:hr} applies with $\ell=p$. By Theorem~\ref{thm:frey-hr},
$\rhobar_{E,p}$ is hardly ramified. Applying $B_5$ with $\ell=p$ to the
two-dimensional $\F_p$-vector space $E(\Qbar)[p]$, we conclude that
$\rhobar_{E,p}$ is reducible. As $(a,b,c,p)$ and $\Qbar$ were arbitrary, this
is exactly Statement $B_4$. Hence $B_5\Rightarrow B_4\Rightarrow{}$FLT.
\end{proof}

\begin{remark}
This reduction is, in essence, the beginning of Serre's 1987 paper: the whole
of the argument so far --- the reduction of FLT to the reducibility of a single
hardly ramified mod $p$ representation --- was available in the 1980s. What was
\emph{not} available then, and what the remaining levels supply, is a proof
that a hardly ramified representation is reducible without ever mentioning the
Frey curve, modular forms, or Serre's conjecture. Statement $B_5$ would follow
from the conjecture Serre stated in that very paper; but Serre's conjecture was
not proved until the twenty-first century, and the argument we shall give uses
ideas from its proof (by Khare and Wintenberger) rather than the conjecture
itself.
\end{remark}

\end{document}
